% !PRINTED 113
\weberpage{113}
%<<CONT

\begin{equation}
 x=\sqrt[n]{r}\left(\cos\frac{\varphi}{n}+i\sin\frac{\varphi}{n}\right)
\tag{7}
\end{equation}
is a magnitude whose $n$th power is $=a$, hence a root of the equation
\begin{equation}
 x^n=a,
\tag{8}
\end{equation}
where $n$ is understood to be an arbitrary positive whole number. But every one of the values
\begin{equation}
 x_k=\sqrt[n]{r}\left[\cos\left(\frac{\varphi+2k\pi}{n}\right)+i\sin\left(\frac{\varphi+2k\pi}{n}\right)\right],
\tag{9}
\end{equation}
satisfies the same requirement, when $k$ is an arbitrary integer. In (9), however, not more than $n$ different values are contained, namely those obtained by putting
\begin{equation}
 k=0,1,2,\ldots,n-1;
\tag{10}
\end{equation}
for if $k$ is increased by a multiple of $n$, the value of (9) is not changed, while the $n$ values (10) give only different values of $x$, because they have different phases.

Equation (8) consequently has not merely one, but $n$ different roots. The geometrical images of the values $x_k$ all lie on a circle with radius $\sqrt[n]{r}$, and are separated from one another by the angle $\frac{2\pi}{n}$. The first of these points is obtained by dividing the given angle $\varphi$ into $n$ equal parts. The radius $\sqrt[n]{r}$ is greater or smaller than $r$ according as $r$ is smaller or greater than $1$. The adjacent figure shows these relations for $n=3$ under the assumption $r>1$.


% Exact 300-dpi crop from the registered Volume-I scan, printed Fig. 1.
\weberfigureimage{weber-b1-fig01.png}{0.50}


One can obtain the different values $x_k$ by multiplying the first of them,
\[
 x_0=\sqrt[n]{r}\left(\cos\frac{\varphi}{n}+i\sin\frac{\varphi}{n}\right),
\]
by the different values of
%CONT>>
